\scanpage{7}
% L01–03
est \unsure{splittée|splittie}, d'où un hom \inserted{$u$} de $H_0$ dans $t_{K_0}$ qui donne l'extension de $\unsure{K_0|K}$ par $\unsure{\mathfrak g K_0|\mathfrak L K_0|\mathbb V t_{K_0}|\mathscr L K_0|J K_0}$ plus l'hom $G\longrightarrow K$ \deleted{\unsure{de}} (donc permet de réconstituer $H$)
% L04
\[ u:H_0\longrightarrow\mathfrak g K_0 \]
% L05–09
\deleted{Ici $H_0=\alpha_p$ \inserted{$\subset G_0=G_a$}, $K_0\simeq G_{a\unsure{p}}$, $\mathfrak g K_0\simeq G_a$, $\operatorname{Hom}(\alpha_p,G_a)$ et on sait $\operatorname{Hom}(\alpha_p,\unsure{\mathrm{fibr\acute e}|\mathrm{faisceau}})\xrightarrow{\sim}\operatorname{Hom}(\mathfrak g\alpha_p,\mathfrak g\unsure{\mathrm{fibr\acute e}|\mathrm{faisceau}})$ donc $\operatorname{Hom}(K_0,\mathfrak g K_0)\mathbin{\text{---}}\operatorname{Hom}(H_0,\mathfrak g K_0)$ surj, donc \deleted{l'extension} \unsure{le prolongt} $\overline K$ de $K_0$ est trivial i.e. $\simeq$}
% L09–10
On trouve l'espace tangent de Zariski :
% L10
\[ \operatorname{Hom}_{\deleted{\unsure{k}}\unsure{k\text{-gr}|\text{gr}}}(H_0,\mathfrak g K_0) \]
% L11
\begin{center}- - -\end{center}

% L12
\uline{Théorème} $\varphi_n:\mathbb E^n_{\mathbb F_p}\longrightarrow X_n$ \uline{iso}

% L13–17
i.e. Si $S$ est un schéma de car $p$ $H\subset G_{a\,S}$, ss-schéma en \unsure{gpes|groupes} fini et \deleted{\unsure{plat}} \inserted{loc. libre} de rang $p^n$, alors $\exists1\ a=(a_0,\text{--},a_{n-1})\in\mathbb E^n(S)=\Gamma(S,\mathscr O_S)^n$ tel que
% L17
\[ H=\operatorname{Ker} f_a \]
% L18
où
% L18
\[ f_a:G_{a\,S}\mathbin{\text{---}}G_{a\,S}\qquad f_a(x)=x^{p^n}+a_{n-1}x^{p^{n-1}}\ \text{--}+a_0x . \]
% L19
De plus, $H$ est étale $/S$ sss $a_0\in\Gamma(S,\mathscr O_S)^*=\Gamma(S,\mathscr O_S^*)$
% archive sheet number: 6
